% GENERATED FILE. Edit canonical lessons, then run npm run generate:books.
% canonical-source-hash: fnv1a64:83f013a10131626b
% canonical-lessons: AR-C421-fursa, AR-C421-laalla, AR-C421-iktishaf, AR-C421-dhaka, AR-C421-ijtihad, AR-R421-review-414-421

\chapter{Chances and discoveries}
\label{ch:ar-chances-and-discoveries}
\hlchaptermodality{421}

\begin{chapteropening}
\textbf{By the end of this chapter:} \emph{I can talk about a chance, a discovery, intelligence and diligence, and say it may be that.}

Complete the last lesson of chapter 421: I can talk about a chance, a discovery, intelligence and diligence, and say it may be that.
\end{chapteropening}

\section[\texorpdfstring{\ar{فرصة} (furṣa)}{فرصة (furṣa)}]{\ar{فرصة} (furṣa) — a chance}

\label{lesson:AR-C421-fursa}



\begin{quote}
Before the new one: say the Arabic for to be equal, then the Arabic for shared.
\end{quote}

\subsection*{You'll want to know: \ar{فرصة}}
\textbf{\ar{فرصة}} — \emph{furṣa} — \textquotedblleft{}a chance\textquotedblright{}.

\begin{quote}
\ar{كانت} \ar{الرحلة} \ar{فرصة} \ar{جيدة} \ar{للتعلم}.
\end{quote}

\begin{quote}
\emph{kānat ar-riḥla furṣa jayyida li-t-taʿallum.}
\end{quote}

\begin{quote}
The trip was a good chance to learn.
\end{quote}

Plural: \ar{فرص}, \emph{furaṣ}.

\begin{cousinweb}
\textbf{\ar{فرصة}} (\emph{furṣa}) is built on the native Arabic root \emph{f-r-ṣ}, \textquotedblleft{}a turn, an opportunity\textquotedblright{} (Lane). A \emph{furṣa} is one's turn, the right moment: a chance.
\end{cousinweb}

\begin{grammarlens}[title={What you've built}]
One more word for saying what happens if.
\end{grammarlens}

\subsection*{Your turn}
\begin{itemize}
\raggedright
  \item \emph{Say it:} \emph{furṣa}
  \item \emph{Say it:} \emph{furṣa}, once more
  \item \emph{Say it:} \emph{kānat ar-riḥla furṣa jayyida li-t-taʿallum.} — The trip was a good chance to learn.
  \item \emph{Recall:} say the Arabic for to be equal, then the Arabic for shared, then say \emph{furṣa} again
\end{itemize}

\begin{tcolorbox}[breakable,colback=teal!4,colframe=teal!35!black,title={Before you move on}]
What does \ar{فرصة} mean? (\textquotedblleft{}A chance\textquotedblright{}.) Say it once more.
\end{tcolorbox}
\section[\texorpdfstring{\ar{لعل} (laʿalla)}{لعل (laʿalla)}]{\ar{لعل} (laʿalla) — it may be that; perhaps}

\label{lesson:AR-C421-laalla}



\begin{quote}
Before the new one: say the Arabic for shared, then the Arabic for a chance.
\end{quote}

\subsection*{You'll want to know: \ar{لعل}}
\textbf{\ar{لعل}} — \emph{laʿalla} — \textquotedblleft{}it may be that; perhaps\textquotedblright{}.

\begin{quote}
\ar{لعل} \ar{الطقس} \ar{يتحسن} \ar{غدا}.
\end{quote}

\begin{quote}
\emph{laʿalla aṭ-ṭaqs yataḥassanu ghadan.}
\end{quote}

\begin{quote}
Perhaps the weather will improve tomorrow.
\end{quote}

Like \ar{إن}, it is followed by a noun or a pronoun: \ar{لعله}, perhaps he.

\begin{cousinweb}
\textbf{\ar{لعل}} (\emph{laʿalla}) is an old native Arabic particle of hope or expectation, perhaps. Grammarians often trace it to \emph{ʿalla} with \emph{la-}; its origin is not settled.
\end{cousinweb}

\begin{grammarlens}[title={What you've built}]
One more word for saying what happens if.
\end{grammarlens}

\subsection*{Your turn}
\begin{itemize}
\raggedright
  \item \emph{Say it:} \emph{laʿalla}
  \item \emph{Say it:} \emph{laʿalla}, once more
  \item \emph{Say it:} \emph{laʿalla aṭ-ṭaqs yataḥassanu ghadan.} — Perhaps the weather will improve tomorrow.
  \item \emph{Recall:} say the Arabic for shared, then the Arabic for a chance, then say \emph{laʿalla} again
\end{itemize}

\begin{tcolorbox}[breakable,colback=teal!4,colframe=teal!35!black,title={Before you move on}]
What does \ar{لعل} mean? (\textquotedblleft{}It may be that; perhaps\textquotedblright{}.) Say it once more.
\end{tcolorbox}
\section[\texorpdfstring{\ar{اكتشاف} (iktishāf)}{اكتشاف (iktishāf)}]{\ar{اكتشاف} (iktishāf) — a discovery}

\label{lesson:AR-C421-iktishaf}



\begin{quote}
Before the new one: say the Arabic for a chance, then the Arabic for it may be that.
\end{quote}

\subsection*{You'll want to know: \ar{اكتشاف}}
\textbf{\ar{اكتشاف}} — \emph{iktishāf} — \textquotedblleft{}a discovery\textquotedblright{}.

\begin{quote}
\ar{كان} \ar{اكتشاف} \ar{الدواء} \ar{مهما} \ar{للعالم}.
\end{quote}

\begin{quote}
\emph{kāna iktishāf ad-dawāʾ muhimman li-l-ʿālam.}
\end{quote}

\begin{quote}
The discovery of the medicine was important for the world.
\end{quote}

\begin{cousinweb}
\textbf{\ar{اكتشاف}} (\emph{iktishāf}) is built on the native Arabic root \emph{k-sh-f}, \textquotedblleft{}to uncover\textquotedblright{} (Lane). \emph{Iktishāf} is the Form VIII verbal noun of \emph{iktashafa}, to discover: a discovery.
\end{cousinweb}

\begin{grammarlens}[title={What you've built}]
One more word for saying what happens if.
\end{grammarlens}

\subsection*{Your turn}
\begin{itemize}
\raggedright
  \item \emph{Say it:} \emph{iktishāf}
  \item \emph{Say it:} \emph{iktishāf}, once more
  \item \emph{Say it:} \emph{kāna iktishāf ad-dawāʾ muhimman li-l-ʿālam.} — The discovery of the medicine was important for the world.
  \item \emph{Recall:} say the Arabic for a chance, then the Arabic for it may be that, then say \emph{iktishāf} again
\end{itemize}

\begin{tcolorbox}[breakable,colback=teal!4,colframe=teal!35!black,title={Before you move on}]
What does \ar{اكتشاف} mean? (\textquotedblleft{}A discovery\textquotedblright{}.) Say it once more.
\end{tcolorbox}
\section[\texorpdfstring{\ar{ذكاء} (dhakāʾ)}{ذكاء (dhakāʾ)}]{\ar{ذكاء} (dhakāʾ) — intelligence}

\label{lesson:AR-C421-dhaka}



\begin{quote}
Before the new one: say the Arabic for it may be that, then the Arabic for a discovery.
\end{quote}

\subsection*{You'll want to know: \ar{ذكاء}}
\textbf{\ar{ذكاء}} — \emph{dhakāʾ} — \textquotedblleft{}intelligence\textquotedblright{}.

\begin{quote}
\ar{الذكاء} \ar{وحده} \ar{لا} \ar{يكفي}.
\end{quote}

\begin{quote}
\emph{adh-dhakāʾ waḥdahu lā yakfī.}
\end{quote}

\begin{quote}
Intelligence alone is not enough.
\end{quote}

\begin{cousinweb}
\textbf{\ar{ذكاء}} (\emph{dhakāʾ}) is built on the native Arabic root \emph{dh-k-w}, usually traced to the sense \textquotedblleft{}to be bright, sharp of mind\textquotedblright{}. \emph{Dhakāʾ} is sharpness of mind: intelligence. \ar{ذكي}, clever, is from the same root.
\end{cousinweb}

\begin{grammarlens}[title={What you've built}]
One more word for saying what happens if.
\end{grammarlens}

\subsection*{Your turn}
\begin{itemize}
\raggedright
  \item \emph{Say it:} \emph{dhakāʾ}
  \item \emph{Say it:} \emph{dhakāʾ}, once more
  \item \emph{Say it:} \emph{adh-dhakāʾ waḥdahu lā yakfī.} — Intelligence alone is not enough.
  \item \emph{Recall:} say the Arabic for it may be that, then the Arabic for a discovery, then say \emph{dhakāʾ} again
\end{itemize}

\begin{tcolorbox}[breakable,colback=teal!4,colframe=teal!35!black,title={Before you move on}]
What does \ar{ذكاء} mean? (\textquotedblleft{}Intelligence\textquotedblright{}.) Say it once more.
\end{tcolorbox}
\section[\texorpdfstring{\ar{اجتهاد} (ijtihād)}{اجتهاد (ijtihād)}]{\ar{اجتهاد} (ijtihād) — diligence}

\label{lesson:AR-C421-ijtihad}



\begin{quote}
Before the new one: say the Arabic for a discovery, then the Arabic for intelligence.
\end{quote}

\subsection*{You'll want to know: \ar{اجتهاد}}
\textbf{\ar{اجتهاد}} — \emph{ijtihād} — \textquotedblleft{}diligence\textquotedblright{}.

\begin{quote}
\ar{نجح} \ar{الطالب} \ar{باجتهاده}.
\end{quote}

\begin{quote}
\emph{najaḥa aṭ-ṭālib bi-jtihādih.}
\end{quote}

\begin{quote}
The student succeeded through hard work.
\end{quote}

\begin{cousinweb}
\textbf{\ar{اجتهاد}} (\emph{ijtihād}) is built on the native Arabic root \emph{j-h-d}, \textquotedblleft{}to strive, to exert oneself\textquotedblright{} (Lane). \emph{Ijtihād} is the Form VIII verbal noun, exerting oneself: diligence. \ar{مجتهد}, hard-working, is from the same root.
\end{cousinweb}

\begin{grammarlens}[title={What you've built}]
One more word for saying what happens if.
\end{grammarlens}

\subsection*{Your turn}
\begin{itemize}
\raggedright
  \item \emph{Say it:} \emph{ijtihād}
  \item \emph{Say it:} \emph{ijtihād}, once more
  \item \emph{Say it:} \emph{najaḥa aṭ-ṭālib bi-jtihādih.} — The student succeeded through hard work.
  \item \emph{Recall:} say the Arabic for a discovery, then the Arabic for intelligence, then say \emph{ijtihād} again
\end{itemize}

\begin{tcolorbox}[breakable,colback=teal!4,colframe=teal!35!black,title={Before you move on}]
What does \ar{اجتهاد} mean? (\textquotedblleft{}Diligence\textquotedblright{}.) Say it once more.
\end{tcolorbox}
\section[(dialogue)]{Words from chapters 414–421 — a review}

\label{lesson:AR-R421-review-414-421}



\begin{quote}
Say the Arabic for to stem from and to follow from.
\end{quote}

\subsection*{Your turn}
Say these aloud:

\begin{itemize}
\raggedright
  \item to stem from, to follow from, to be reflected, a consequence, an effect
  \item an achievement, to be affected, to become clear, to arrive at, a level
  \item to adopt, instead of, in addition to, insurance, apart from
  \item vague, frank, specific, to become evident, to refer
  \item to strengthen, to weaken, to develop, effectiveness, independent
  \item positive, negative, basic, secondary, comprehensive
  \item to respond, to ignore, to nominate, to be equal, shared
  \item a chance, it may be that, a discovery, intelligence, diligence
\end{itemize}

\begin{tcolorbox}[breakable,colback=teal!4,colframe=teal!35!black,title={Before you move on}]
To stem from? (\textbf{\ar{نتج}}, \emph{nataja}.) To follow from? (\textbf{\ar{ترتب}}, \emph{tarattaba}.) To be reflected? (\textbf{\ar{انعكس}}, \emph{inʿakasa}.) A consequence? (\textbf{\ar{عاقبة}}, \emph{ʿāqiba}.) An effect? (\textbf{\ar{تأثير}}, \emph{taʾthīr}.) An achievement? (\textbf{\ar{إنجاز}}, \emph{injāz}.) To be affected? (\textbf{\ar{تأثر}}, \emph{taʾaththara}.) To become clear? (\textbf{\ar{تبين}}, \emph{tabayyana}.) To arrive at? (\textbf{\ar{توصل}}, \emph{tawaṣṣala}.) A level? (\textbf{\ar{مستوى}}, \emph{mustawā}.) To adopt? (\textbf{\ar{تبنى}}, \emph{tabannā}.) Instead of? (\textbf{\ar{بدلا} \ar{من}}, \emph{badalan min}.) In addition to? (\textbf{\ar{بالإضافة} \ar{إلى}}, \emph{bi-l-iḍāfa ilā}.) Insurance? (\textbf{\ar{تأمين}}, \emph{taʾmīn}.) Apart from? (\textbf{\ar{سوى}}, \emph{siwā}.) Vague? (\textbf{\ar{غامض}}, \emph{ghāmiḍ}.) Frank? (\textbf{\ar{صريح}}, \emph{ṣarīḥ}.) Specific? (\textbf{\ar{محدد}}, \emph{muḥaddad}.) To become evident? (\textbf{\ar{اتضح}}, \emph{ittaḍaḥa}.) To refer? (\textbf{\ar{أشار}}, \emph{ashāra}.) To strengthen? (\textbf{\ar{عزز}}, \emph{ʿazzaza}.) To weaken? (\textbf{\ar{أضعف}}, \emph{aḍʿafa}.) To develop? (\textbf{\ar{طور}}, \emph{ṭawwara}.) Effectiveness? (\textbf{\ar{فعالية}}, \emph{faʿʿāliyya}.) Independent? (\textbf{\ar{مستقل}}, \emph{mustaqill}.) Positive? (\textbf{\ar{إيجابي}}, \emph{ījābī}.) Negative? (\textbf{\ar{سلبي}}, \emph{salbī}.) Basic? (\textbf{\ar{أساسي}}, \emph{asāsī}.) Secondary? (\textbf{\ar{ثانوي}}, \emph{thānawī}.) Comprehensive? (\textbf{\ar{شامل}}, \emph{shāmil}.) To respond? (\textbf{\ar{استجاب}}, \emph{istajāba}.) To ignore? (\textbf{\ar{تجاهل}}, \emph{tajāhala}.) To nominate? (\textbf{\ar{رشح}}, \emph{rashshaḥa}.) To be equal? (\textbf{\ar{تساوى}}, \emph{tasāwā}.) Shared? (\textbf{\ar{مشترك}}, \emph{mushtarak}.) A chance? (\textbf{\ar{فرصة}}, \emph{furṣa}.) It may be that? (\textbf{\ar{لعل}}, \emph{laʿalla}.) A discovery? (\textbf{\ar{اكتشاف}}, \emph{iktishāf}.) Intelligence? (\textbf{\ar{ذكاء}}, \emph{dhakāʾ}.) Diligence? (\textbf{\ar{اجتهاد}}, \emph{ijtihād}.)
\end{tcolorbox}
